\subsection*{Parity signs and certified evaluation of the inverse refinement}
The global remainder norm in the preceding inverse estimate loses useful
signs. Retaining them improves the complete far-operator sandwich;
it does not determine the sign of the low Schur matrix.

\begin{proposition}[Sharper inverse constants at the fixed window]
\label{prop:v122-parity-inverse}
For the actual shifted far operators and the unchanged weights $D_g$
in Proposition~\ref{prop:v120-lambda4-tail}, one has
\[
 D_g\preceq U\preceq20D_g\quad\text{in the even sector }n>512,
 \qquad
 D_g\preceq U\preceq90D_g\quad\text{in the odd sector }n>1536.
\]
Thus the constants in Proposition~\ref{prop:v120-inverse-polynomial}
can be replaced by $20$ and $90$, respectively, with error factors
$(19/20)^k$ and $(89/90)^k$. These are complete-operator bounds
at $\lambda=4$, with no assertion of growing-window uniformity.
\end{proposition}
\begin{proof}
Put $h=9/16$, $c=10^{-8}$, $L=2\log4$ and
$t_N=2\pi(N+1)/L$. Let $M_\phi$ be the certified prime bound
in \eqref{eq:v120-lambda4-prime-bound}. The limiting archimedean
off-diagonal in the normalized parity basis is the matrix
$1/[2(n+m)]$ in the even sector and its negative in the odd sector.
It is positive and has norm at most $\pi/2$; the remaining
archimedean commutator has norm at most $2/t_N$, by the proof of
Proposition~\ref{prop:v115-sharp-tail}.

In centered coordinates the pole form is
$2|\ip f{\cosh(x/2)}|^2-2|\ip f{\sinh(x/2)}|^2$.
Only the positive term acts on even vectors, and only the negative
term on odd vectors. Direct integration gives
\[
 |\widehat{\cosh(x/2)}(n)|^2
 =\frac{h}{L(t_n^2+1/4)^2},\qquad
 2\|Q_N\cosh(x/2)\|^2
 \le\frac{hL^3}{4\pi^4}\sum_{n>N}n^{-4}
 \le\frac{hL^3}{12\pi^4N^3}.
\]
The factor two for the form and both Fourier signs are included.
Consequently the remainder above the shifted archimedean diagonal
has the following upper bounds:
\begin{equation}\label{eq:v122-parity-upper}
 \begin{split}
 C_{\rm e}(N)&=M_\phi+2/t_N+\pi/2+hL^3/(12\pi^4N^3),\\
 C_{\rm o}(N)&=M_\phi+2/t_N.
 \end{split}
\end{equation}
The odd limiting Hilbert matrix and odd pole term are nonpositive
and may be omitted in an upper bound. The shift $cD$ changes only
the exact archimedean diagonal; none of these remainder terms is
multiplied by $1-c$.

For the same $\kappa$ and increasing $g_n$ as before, the two-sided
diagonal estimate \eqref{eq:v120-two-sided-arch} now yields
\[
 U\preceq D_g+
 \{\kappa+2(1-c)E_L(t_N)+C_{\rm e/o}(N)\}I.
\]
It suffices to bound one plus the braced quantity divided by
$g_{N+1}$. Outward interval evaluation gives values below
$19.215607$ for even $N=512$, and below $89.840960$ for odd
$N=1536$. The integer gates $20$ and $90$ pass at both 320 and
384 bits. The lower bounds and closed form domains are unchanged.
\end{proof}

\begin{proposition}[Complete residual bounds from a finite prefix]
\label{prop:v122-prefix-refinement}
Assume $D_g\preceq U\preceq mD_g$, with
$D_g\succeq\gamma I>0$. Put
$A=D_g^{-1/2}UD_g^{-1/2}$, $H=A-I$, and let $P$ be an
orthogonal prefix projection, $Q=I-P$. Suppose
$QHQ\preceq\nu Q$ with $\nu\ge0$. For a residual $k$ write
\[
 y=D_g^{-1/2}k=p+q,\quad p=Py,\quad q=Qy,\qquad
 v=\norm p^2,\quad a=\ip{Hp}{p},\quad \norm q^2\le E.
\]
With $[s]_+=\max(s,0)$, the following scalar bound includes
every omitted residual row:
\begin{equation}\label{eq:v122-prefix-scalar}
 \ip{U^{-1}k}{k}
 \le v+E-\frac1m[\sqrt a-\sqrt{\nu E}]_+^2.
\end{equation}
For a matrix of residual columns $Y=D_g^{-1/2}\mathcal R$, put
$Y_P=PY$, $Y_Q=QY$, $V=Y_P^*Y_P$ and
$A_P=Y_P^*HY_P$. If $Y_Q^*Y_Q\preceq\mathcal E$, then, for
every fixed $0<\theta\le1$,
\begin{equation}\label{eq:v122-prefix-matrix}
 \mathcal R^*U^{-1}\mathcal R\preceq
 V+\mathcal E-\frac1m
 \{(1-\theta)A_P-(\theta^{-1}-1)\nu\mathcal E\}.
\end{equation}
If a certified $0\le\rho<1$ instead satisfies
$\nu\mathcal E\preceq\rho^2A_P$, the braced lower bound can
be replaced by $(1-\rho)^2A_P$.
\end{proposition}
\begin{proof}
The bounded operator $H$ is positive, so the reverse triangle
inequality for $H^{1/2}p+H^{1/2}q$ gives
\[
 \ip{Hy}{y}\ge
 [\sqrt a-\|H^{1/2}q\|]_+^2
 \ge[\sqrt a-\sqrt{\nu E}]_+^2.
\]
Combine this with the first inverse polynomial
$A^{-1}\preceq I-H/m$ and
$\norm y^2=\norm p^2+\norm q^2\le v+E$.
For the matrix version, apply the lower Young inequality to
$H^{1/2}Y_Px+H^{1/2}Y_Qx$ for every head vector $x$.
Since $1-\theta^{-1}\le0$, the remote upper bound supplies the
stated lower form. The relative variant follows from the scalar
reverse triangle inequality in every head direction. No factor
equal to the number of columns is inserted. The scalar positive
part in \eqref{eq:v122-prefix-scalar} is not a matrix positive-part
operation; the latter does not follow from these arguments.
\end{proof}

In the actual parity tail based at $N$, the prefix can end at any
$J>N$. The same proof as \eqref{eq:v122-parity-upper} gives the
complete remote compression bound
\begin{equation}\label{eq:v122-remote-H}
 Q_JHQ_J\preceq\nu_J Q_J,\qquad
 \nu_J=\frac{\kappa_N+2(1-c)E_L(t_J)+C_{\rm e/o}(J)}{g_{J+1}}.
\end{equation}
Here $t_J=2\pi(J+1)/L$, whereas the constant $\kappa_N$ in
the original $D_g$ is kept fixed. The upper numerator decreases
and the positive denominator increases with the index. The moment
bound of Proposition~\ref{prop:v116-remote-moment}, divided by
$g_{J+1}$, supplies $\mathcal E$. Thus these formulas evaluate a
signed inverse improvement without setting the infinite residual
to zero. All occurrences of $U$ are interpreted through its
bounded preconditioned form.

\begin{proposition}[Refining both terms of the structured inverse]
\label{prop:v122-refined-structured}
Keep the factors $Z,R,C,\underline L,\mu$ of
Proposition~\ref{prop:v120-structured-inverse}, and assume also
$U\preceq mD_g$. Let $Q_*$ be a bounded positive operator such
that
\[
 U^{-1}\preceq Q_*\preceq D_g^{-1}.
\]
Then, for every $r=(h,k)$,
\begin{equation}\label{eq:v122-refined-structured}
 \ip{T^{-1}r}{r}
 \le\ip{Q_*k}{k}
       +\mu^{-1}\norm{C(h-Z^*k-R^*Q_*k)}^2.
\end{equation}
The same inequality holds for the simultaneous residual Gram.
In particular all inverse polynomials of
Proposition~\ref{prop:v120-inverse-polynomial} can be used while
retaining the original certified $C$ and $\mu$.
\end{proposition}
\begin{proof}
The assumed sandwich implies
$m^{-1}D_g^{-1}\preceq Q_*\preceq D_g^{-1}$.
Consequently $Q_*$ is injective and the closed form of $Q_*^{-1}$
has the same domain as $D_g^{1/2}$, with equivalent form norms.
Inverse order gives $U\succeq Q_*^{-1}$ as forms. Moreover
\[
 K_Z-R^*Q_*R\succeq K_Z-R^*D_g^{-1}R\succeq\underline L.
\]
Repeat the lower-form square completion in
Proposition~\ref{prop:v120-structured-inverse}, replacing the far
form by $Q_*^{-1}$ and its inverse by $Q_*$. Its exact inverse
contains $Q_*$ in both displayed terms, proving the assertion.
The change preserves the lower form domain since $Q_*Rx$ belongs
to the operator domain of $Q_*^{-1}$ for every finite head vector.
Applying the result to all linear combinations proves the Gram
version. If the saved factors certify $T-cD$, use inverse order
as before and retain those shifted factors throughout.
\end{proof}

Improving only the first term while leaving the old mixed term
unchanged is not justified. Positivity of the complete low Schur
matrix requires the resulting full correction, on the entire head,
to fit below $K_X$. The new estimates do not assert that comparison
or convergence along growing windows.

\begin{remark}[The simultaneous head acceptance test]
\label{rem:v122-head-ldl}
For a positive definite trial matrix $K_X$, put
$C_X=r^*T^{-1}r$. Then
\[
 \lambda_{\max}(K_X^{-1/2}C_XK_X^{-1/2})\le1
 \quad\Longleftrightarrow\quad K_X-C_X\succeq0.
\]
Thus an eigenvalue or inverse-square-root estimate is unnecessary.
It suffices to construct a Hermitian \emph{upper enclosure in form
order} $\overline C_X\succeq C_X$ for the complete residual Gram
and certify $K_X-\overline C_X\succeq0$, for example by interval
$LDL^*$ with strictly positive pivots. A fixed invertible dyadic
congruence may improve conditioning without changing the sign.
Strictly positive pivots prove the stronger strict comparison;
an exact zero pivot requires the corresponding semidefinite
factorization conditions, not an interval containing zero.
An entrywise upper bound is not by itself an upper bound in form
order. For $T_M=P_MT P_M$ and $r_M=P_Mr$, the variational identity
gives $r_M^*T_M^{-1}r_M\preceq C_X$, the opposite enclosure
direction. All mixed and remote contributions must be enclosed
before the $LDL^*$ test proves the complete Schur sign.
\end{remark}
